% Options for packages loaded elsewhere
\PassOptionsToPackage{unicode}{hyperref}
\PassOptionsToPackage{hyphens}{url}
\documentclass[
  ignorenonframetext,
]{beamer}
\newif\ifbibliography
\usepackage{pgfpages}
\setbeamertemplate{caption}[numbered]
\setbeamertemplate{caption label separator}{: }
\setbeamercolor{caption name}{fg=normal text.fg}
\beamertemplatenavigationsymbolsempty
% remove section numbering
\setbeamertemplate{part page}{
  \centering
  \begin{beamercolorbox}[sep=16pt,center]{part title}
    \usebeamerfont{part title}\insertpart\par
  \end{beamercolorbox}
}
\setbeamertemplate{section page}{
  \centering
  \begin{beamercolorbox}[sep=12pt,center]{section title}
    \usebeamerfont{section title}\insertsection\par
  \end{beamercolorbox}
}
\setbeamertemplate{subsection page}{
  \centering
  \begin{beamercolorbox}[sep=8pt,center]{subsection title}
    \usebeamerfont{subsection title}\insertsubsection\par
  \end{beamercolorbox}
}
% Prevent slide breaks in the middle of a paragraph
\widowpenalties 1 10000
\raggedbottom
\AtBeginPart{
  \frame{\partpage}
}
\AtBeginSection{
  \ifbibliography
  \else
    \frame{\sectionpage}
  \fi
}
\AtBeginSubsection{
  \frame{\subsectionpage}
}
\usepackage{iftex}
\ifPDFTeX
  \usepackage[T1]{fontenc}
  \usepackage[utf8]{inputenc}
  \usepackage{textcomp} % provide euro and other symbols
\else % if luatex or xetex
  \usepackage{unicode-math} % this also loads fontspec
  \defaultfontfeatures{Scale=MatchLowercase}
  \defaultfontfeatures[\rmfamily]{Ligatures=TeX,Scale=1}
\fi
\usepackage{lmodern}
\ifPDFTeX\else
  % xetex/luatex font selection
\fi
% Use upquote if available, for straight quotes in verbatim environments
\IfFileExists{upquote.sty}{\usepackage{upquote}}{}
\IfFileExists{microtype.sty}{% use microtype if available
  \usepackage[]{microtype}
  \UseMicrotypeSet[protrusion]{basicmath} % disable protrusion for tt fonts
}{}
\makeatletter
\@ifundefined{KOMAClassName}{% if non-KOMA class
  \IfFileExists{parskip.sty}{%
    \usepackage{parskip}
  }{% else
    \setlength{\parindent}{0pt}
    \setlength{\parskip}{6pt plus 2pt minus 1pt}}
}{% if KOMA class
  \KOMAoptions{parskip=half}}
\makeatother
\usepackage{longtable,booktabs,array}
\newcounter{none} % for unnumbered tables
\usepackage{calc} % for calculating minipage widths
\usepackage{caption}
% Make caption package work with longtable
\makeatletter
\def\fnum@table{\tablename~\thetable}
\makeatother
\setlength{\emergencystretch}{3em} % prevent overfull lines
\providecommand{\tightlist}{%
  \setlength{\itemsep}{0pt}\setlength{\parskip}{0pt}}
\usepackage{bookmark}
\IfFileExists{xurl.sty}{\usepackage{xurl}}{} % add URL line breaks if available
\urlstyle{same}
\hypersetup{
  hidelinks,
  pdfcreator={LaTeX via pandoc}}

\author{\texorpdfstring{}{}}
\date{}

\begin{document}

\begin{frame}[fragile]{Scalability: From 1 to N Projects}
\protect\phantomsection\label{scalability-from-1-to-n-projects}
The Standalone Project Paradigm enables horizontal scaling: adding a new
project requires creating a directory with
\texttt{manuscript/config.yaml} and nothing else. No infrastructure code
changes, no \texttt{pyproject.toml} modifications, no CI configuration
updates. The \texttt{run.sh} orchestrator automatically discovers new
projects and presents them in its interactive menu.

We have validated this scaling model with three heterogeneous projects:

\begin{itemize}
\tightlist
\item
  \textbf{\texttt{code\_project}}: Numerical optimization example paper
  with gradient descent, 39 tests, 90\%+ coverage. Demonstrates the
  minimal viable project footprint: a single \texttt{src/} module, a
  single script, and a compact manuscript.
\item
  \textbf{\texttt{act\_inf\_metaanalysis}}: Active inference
  meta-analysis pipeline, 505 tests, 90\%+ coverage, spanning evidence
  synthesis, bibliometric analysis, and citation-weighted knowledge
  graphs. Demonstrates the template's capacity for computationally
  intensive, data-heavy research workflows with large test suites.
\item
  \textbf{\texttt{template}}: This self-referential architectural
  analysis, 65 tests, 90\%+ coverage. Demonstrates the system's ability
  to analyze and document itself---a unique stress test of the Two-Layer
  Architecture's reflexive capability.
\end{itemize}

These projects share no code with each other. They share only the
infrastructure layer---12 modules, \textasciitilde150 Python
files---which provides logging, rendering, validation, steganography,
and reporting services identically to each project. This validates the
Two-Layer Architecture's claim that infrastructure and project concerns
can be cleanly separated.

\begin{block}{Multi-Project Orchestration}
\protect\phantomsection\label{multi-project-orchestration}
When the \texttt{-\/-all-projects} flag is passed to \texttt{run.sh},
the pipeline executes each discovered project sequentially, running
infrastructure tests once at the start and skipping them for individual
projects to avoid redundant validation. After all projects complete, a
cross-project executive report aggregates metrics (test counts, coverage
percentages, page counts, rendering durations) into a unified dashboard
with both JSON and Markdown output formats. This executive reporting
stage provides repository-level visibility without requiring any
project-specific reporting code.
\end{block}

\begin{block}{Scaling Metrics}
\protect\phantomsection\label{scaling-metrics}
{\def\LTcaptype{none} % do not increment counter
\begin{longtable}[]{@{}lccc@{}}
\toprule\noalign{}
Metric & \texttt{code\_project} & \texttt{act\_inf\_metaanalysis} &
\texttt{template} \\
\midrule\noalign{}
\endhead
Source modules & 1 & 12+ & 5 \\
Test files & 1 & 9 & 4 \\
Test count & 39 & 505 & 65 \\
Manuscript chapters & 8 & 14 & 18 \\
Analysis scripts & 1 & 3 & 2 \\
Figures (auto-generated) & 3 & 10+ & 4 \\
\bottomrule\noalign{}
\end{longtable}
}

The infrastructure overhead per project is constant regardless of
project size: the same 12 modules, the same 9 pipeline stages, the same
rendering and validation logic. This O(1) infrastructure cost is the
architectural payoff of the Two-Layer separation.
\end{block}
\end{frame}

\end{document}
